%\pagestyle{mystyle}
\chapter[Adapting Entity Extraction Techniques to the Legal Domain][Adapting EE to Legal Domain]{Adapting Entity Extraction Techniques to the Legal Domain}
\label{chap:extraction}


This chapter presents the contributions for \ref{ob:extraction_assessment}---through an assessment of existing \ac{ee} techniques applied to the legal domain.
First, in \Cref{sec:ijckg}, we address the \emph{novel entity challenge}---an important problem for the legal domain, since court and investigative documents often mention persons that are not known in public \acl{kr}, but whose recognition and linking may enable tasks like precedent case retrieval or interlinking heterogeneous documents mentioning suspects during an investigation.

Secondly we study to which extent general-domain \ac{nlp} approaches for \acl{ee} work on the Italian legal domain.
Initially focusing on Italian judgments from civil courts in \Cref{sec:eejudgements} and in \Cref{sec:adaptation}, where we evaluate different adaptation strategies for \acl{ner}. Then, in \Cref{sec:eechats} we focus on chat logs from an Italian investigation. Indeed, as anticipated in \Cref{sec:intro}, annotating legal documents with \ac{ee} systems may be useful for simplifying the analysis of vast amount of heterogeneous documents (challenges \ref{ch:scalability} and \ref{ch:heterogeneity})---for example, by implementing faceted search applications~\cite{ARMENTANO2014_nlp_faceted,hearst_2006_faceted}.
Finally, \Cref{sec:ee:conclusion} concludes this chapter.

The main contributions of this chapter can be summarized as follows:
\begin{enumerate}
    \item introduction of the incremental \ac{nel} task, with a methodology to adapt static \ac{nel} datasets and the release of a benchmark and baseline pipeline;
    \item evaluation of \ac{ner} and incremental \ac{nel} on Italian civil judgments, including a study of domain adaptation strategies;
    \item creation of two annotated benchmarks for Italian civil judgments and investigative \ac{ima} data, and assessment of an \ac{ee} pipeline on both settings.
\end{enumerate}



The publications related to this chapter are \textcite{ijckg2022,BELLANDI2024,aixia_2023,pozzi_wise_2025}.

